\documentclass[11pt,reqno]{amsart}
\usepackage[T1]{fontenc}
\usepackage{amsmath,amssymb,booktabs,tabularx}
\usepackage[hidelinks]{hyperref}
\usepackage{microtype}
\setlength{\emergencystretch}{1em}
\newcommand{\F}{\mathbb F}
\newcommand{\Q}{\mathbb Q}
\newcommand{\Z}{\mathbb Z}
\newcommand{\Brent}{\mathcal V}
\newcommand{\GL}{\operatorname{GL}}
\newcommand{\rank}{\operatorname{rank}}
\newcommand{\tr}{\operatorname{tr}}
\newtheorem{theorem}{Theorem}[section]
\newtheorem{lemma}[theorem]{Lemma}
\newtheorem{proposition}[theorem]{Proposition}
\title[Local geometry of 47- and 48-term schemes]{Exact local and arithmetic geometry of 47- and 48-term schemes for \(4\times4\) matrix multiplication}
\author{Edson Soares de Souza}
\date{23 November 2026. Version 1.1.0.}
\subjclass[2020]{Primary 68W30, 14B10; Secondary 15A69, 13H05}
\keywords{Brent equations, tensor rank, matrix multiplication, characteristic two, Hensel lifting, rational deformations}
\begin{document}
\begin{abstract}
We certify that five inequivalent 47-term schemes are smooth points of the ordered Brent scheme in characteristic two, of local dimension \(139\), with open symmetry orbits. Gamma47-B is equivalent to C5. Exact first-lift obstructions imply that the integral Brent scheme is supported in characteristic two near these points, excluding ramified mixed-characteristic specializations. For the rational Dumas--Pernet--Sedoglavic 48-term scheme, we give an explicit rational equivalence with its more accurate variant. We also construct rational curves that vary twelve summands of a published one-parameter scheme and transport them along that family. This gives four exact rational surfaces and five independent transverse tangent directions. At \(D(2)\) over \(\F_3\), a sixth direction requires at least sixteen external summands, and the bound is attained. We also exclude zero summands and two-factor collisions in the affine symmetry-saturated closure of each surface. The 48-term theorems concern specific deformations and their affine closures.
\end{abstract}
\maketitle
\markboth{Local geometry of 47- and 48-term schemes}{Local geometry of 47- and 48-term schemes}

\section{Introduction}
A short matrix multiplication scheme can serve as a starting point for two different searches. One can try to lift a binary scheme to characteristic zero without increasing the number of products. One can also deform a rational scheme until a pair of summands merges or one summand disappears. A search result for either operation depends on the particular point or family being searched. We use the ordered Brent scheme to certify when these operations are excluded.

The \(47\)-term part concerns binary schemes and their characteristic-zero lifts. Five representatives have tangent space equal to the infinitesimal symmetry orbit. A parity certificate then detects a stronger arithmetic obstruction, namely that the local integral scheme already lies in characteristic two. The \(48\)-term part starts from the rational construction of Dumas, Pernet and Sedoglavic~\cite{dps}. We construct four rational surfaces through it and control their affine limits under basis changes. A separate calculation at the reduction of \(D(2)\) counts how many external summands a tangent must move to go beyond the five directions supplied by these surfaces. These results restrict particular searches for \(47\) rational products. The existence of a \(47\)-term rational scheme is a global question.

Our companion paper proves an exact local dimension at the recursive \(49\)-term Strassen-square point~\cite{strassensquare}. That point has a tangent space larger than its local dimension, so quadratic obstructions are essential to the dimension proof. Here the \(47\)-term local dimension follows from the orbit, and the \(48\)-term results concern rational families and support restrictions. The \(49\)-term numerical study of Li, Wang and Hu~\cite{lwh} has its own choice of points.

\section{Brent equations and group action}
Let \(k\) be a field and let \(E=\operatorname{Mat}_{4\times4}(k)\). We use two equivalent conventions. In direct output coordinates the matrix multiplication tensor is
\[
 T=\sum_{i,j,k=1}^4 e_{ik}\otimes e_{kj}\otimes e_{ij}.
\]
We write the \(s\)-th ordered summand as \((A_s,B_s,C_s)\) in cyclic matrix coordinates, where \(C_s\) is the transpose of its direct output factor. The integral ordered Brent scheme \(\Brent_r\) is defined by the \(4,096\) scalar equations
\begin{equation}\label{eq:brent}
 F_r(A,B,C)=\sum_{s=1}^r A_s\otimes B_s\otimes C_s^\top-T=0
\end{equation}
in \(48r\) variables. We work with this equation scheme when taking tangent spaces and local rings.

The connected acting group has three copies of \(\GL_4\) and two rescaling parameters per term. Write a group element as \((G_0,G_1,G_2,(\lambda_{s,m}))\) with \(\prod_m\lambda_{s,m}=1\) for each \(s\). Its action in cyclic notation takes the mode-\(m\) factor \(X_{s,m}\) to
\begin{equation}\label{eq:action}
 \lambda_{s,m}G_mX_{s,m}G_{m+1}^{-1},
\end{equation}
with vertices taken modulo three. A simultaneous scalar in each \(G_m\) can be canceled by the term rescalings, giving a three-dimensional kernel. We also use term permutations and permutations of the three modes. A cyclic permutation acts directly on cyclic factors and a reversal transposes each matrix. Below \(J_p=d F_r(p)\), \(K_p=\ker J_p\), and \(O_p\) is the image of the differential of the connected action at the identity.

Differentiating \eqref{eq:brent}, the column for changing the \(i\)-th coordinate of \(A_s\) is \(e_i\otimes B_s\otimes C_s^{\top}\). The corresponding columns for the other two modes are obtained by replacing their factor by the coordinate unit. Thus each column is determined by the two unchanged factors. The rank certificate processes the columns in summand, mode, and row-major coordinate order. It reduces a column by already accepted pivots, normalizes its first remaining nonzero entry, and appends it exactly when its remainder is nonzero. The binary calculation needs only additions. This gives the following dimensions for the displayed points.
\begin{center}
\small
\begin{tabular}{@{}lcrrr@{}}
\toprule
Point or tables & Field & \(\rank J_p\) & \(\dim O_p\) & \(\dim K_p/O_p\)\\
\midrule
Five \(47\)-term representatives & \(\F_2\) & \(2117\) & \(139\) & \(0\)\\
DPS48 and its accurate variant & \(\Q,\F_3\) & \(2153\) & \(141\) & \(10\)\\
\(D(2)\) & \(\F_3\) & \(2153\) & \(141\) & \(10\)\\
\bottomrule
\end{tabular}
\end{center}
For the rational Jacobian rank, the nonzero modular pivots give the lower bound. The exact rational kernel certificate gives the matching upper bound by exhibiting \(151\) independent kernel vectors. The replay checks annihilation by integer multiplication and independence by reduction modulo a prime.

\section{Five smooth orbits in characteristic two}
\begin{theorem}\label{thm:rigid}
At each of the five binary schemes AT47, KM47, C3, C4, and C5, the ordered Brent scheme over \(\overline{\F}_2\) is smooth of local dimension \(139\). Its connected symmetry orbit contains a Zariski-open neighborhood of the point. The five points are pairwise inequivalent under all the actions above. Gamma47-B is equivalent to C5.
\end{theorem}
\begin{proof}
For each table direct tensor contraction gives \eqref{eq:brent}, and the column reduction above gives rank \(2117\) in \(2256\) variables. We derive the orbit rank by a smaller certificate. At vertex \(m\), form the product \(L_{s,m}=X_{s,m}X_{s,m+1}X_{s,m+2}\). It transforms by conjugation by \(G_m\), because the term scalings cancel. An infinitesimal stabilizer with vertex generator \(Y_m\) therefore satisfies \([Y_m,L_{s,m}]=0\). The table lists loop indices whose stacked commutator equations in the sixteen entries of \(Y_m\) have rank \(15\) over \(\F_2\).
\begin{center}
\small
\begin{tabular}{@{}lccc@{}}
\toprule
Scheme & Vertex \(0\) & Vertex \(1\) & Vertex \(2\)\\
\midrule
AT47 & \(2,15,12\) & \(2,3,4,14\) & \(2,15,6\)\\
KM47 & \(2,10,8\) & \(18,6,1\) & \(2,10,4\)\\
C3 & \(2,11,1\) & \(2,1,3\) & \(2,8,1\)\\
C4 & \(4,1,9\) & \(4,6,1\) & \(4,1,2\)\\
C5 & \(11,47\) & \(11,47\) & \(11,6,3\)\\
\bottomrule
\end{tabular}
\end{center}
These are one-based summand indices. A row of the small linear system is obtained by taking one coordinate of \([E_{ab},L_{s,m}]\) for each matrix unit \(E_{ab}\). Thus each \(Y_m\) is scalar, and the vanishing of the variation of each nonzero factor forces its infinitesimal term scaling to cancel the difference of the adjacent vertex scalars. The differential kernel has dimension three, giving orbit rank \(3\cdot16+2\cdot47-3=139\). The tangent-space upper bound on local dimension and the lower bound from the dimension of the orbit image therefore match. A finite type scheme over the algebraic closure with local dimension equal to tangent dimension is regular at this point and thus smooth. Its local irreducible component is unique. The orbit is locally closed and has full dimension in that component, so contains an open neighborhood of the point.

For each summand, sort the three individual matrix ranks. The resulting multiset is invariant under \eqref{eq:action}, term permutations and mode permutations. The triples \((1,1,1),(1,1,3),(1,2,2)\) already distinguish the five representatives. Their multiplicities in AT47, KM47, C3, C4 and C5, respectively, are
\[
 (3,7,0),\qquad (2,8,2),\qquad (2,6,5),\qquad (1,9,0),\qquad (1,9,6).
\] For Gamma47-B, cyclically reorder its modes as \((2,0,1)\), using zero-based indices. Apply \eqref{eq:action} with matrices
\[
 \begin{pmatrix}1&1&1&1\\0&1&0&0\\0&0&1&0\\0&0&0&1\end{pmatrix},\
 \begin{pmatrix}1&0&0&0\\1&1&0&0\\1&0&1&0\\1&0&0&1\end{pmatrix},\
 \begin{pmatrix}1&1&1&0\\0&1&1&1\\1&0&1&1\\1&1&0&1\end{pmatrix}.
\]
The supplied permutation matches all \(47\) transformed terms exactly.
\end{proof}

\section{Mixed-characteristic specializations}
\begin{theorem}\label{thm:ramified}
At each point of Theorem~\ref{thm:rigid}, the local ring of the integral Brent scheme over \(\Z_{(2)}\) is a regular local ring of characteristic two. Consequently no mixed-characteristic discrete valuation ring, ramified or unramified, supports a Brent solution specializing to such a point, even after extending the residue field.
\end{theorem}
\begin{proof}
Let \(p\) denote the binary integer representative. Its integer residual is divisible by two. The mod-four equation for \(p+2x\) is
\[
 J_p x= - F_{47}(p)/2\pmod 2.
\]
The checker uses two independent descriptions of the obstruction. It row-reduces the residual in the \(1,979\)-dimensional cokernel of \(J_p\), and verifies a binary functional \(\ell\) with \(\ell J_p=0\), \(\ell(F_{47}(p)/2)=1\). The appendix records the projection weights. The supports of the functionals are frozen in the artifact.

Let \(A\) be the local ring of \(\Brent_{47}\) over \(\Z_{(2)}\) at the binary point. The ambient cotangent space has \(2,257\) generators, one for each coordinate difference and one for \(2\). The linear classes of the equations in this cotangent space are the rows of \([J_p\mid F_{47}(p)/2]\). This matrix has rank \(2,118\) by the functional above, so the embedding dimension of \(A\) equals \(139\). The special-fiber local ring \(A/(2)\) has dimension \(139\) by Theorem~\ref{thm:rigid}. Dimension bounds give \(\dim A=139\), hence \(A\) is regular. A regular local ring is a domain. If \(2\) were nonzero in \(A\), its quotient by \(2\) would have strictly smaller dimension. Thus \(2=0\) in \(A\). The same equality holds on a neighborhood, by the definition of localization. A mixed-characteristic specialization through the specified point would factor locally through this neighborhood.
\end{proof}

\section{Rational 48-term schemes}
At the DPS48 base point of Dumas, Pernet and Sedoglavic~\cite{dps} and its accurate variant, the exact ranks over \(\Q\) and \(\F_3\) give \(\dim K=151\), \(\dim O=141\), and transverse dimension \(10\). All \(55\) quadratic coefficient classes vanish, and the Cycle 90 construction simultaneously lifts the ten directions through degree four over \(\F_3\). Finite-order data in characteristic three imposes no uniqueness of rational coefficients or of higher-order terms.

\begin{proposition}\label{prop:accurate}
The two rational DPS tables are equivalent over \(\Q\) under \eqref{eq:action}, cyclic mode rotation, and term permutations.
\end{proposition}
\begin{proof}
The certificate gives nonsingular integer vertex matrices \(G_m\), rational layer scalings, and a term permutation. The verifier checks the change of basis without explicitly forming inverses by multiplying \eqref{eq:action} on the right by \(G_{m+1}\). It checks all \(48\) term equalities and the product-one constraint on the three layer scalings.
\end{proof}

\subsection*{Boundary of the published one-parameter family}
Let \(D(t)\) be the Li--Wang--Hu family on \(\mathbb G_m\)~\cite{lwh}. At \(t=1/4\) it equals DPS48 up to layer rescaling and term order.
\begin{theorem}\label{thm:lwh}
Over \(\overline{\Q}\) and \(\overline{\F}_3\), the affine closure of the symmetry saturation of \(D(t)\) has no zero summand and no pair of summands with proportional factors in two modes.
\end{theorem}
\begin{proof}
At each \(t\ne0,\infty\), each trace \(\tr(A_s B_s C_s)\) is a unit \(1\) or \(2\). For every pair and mode, a verified coordinate of the wedge product is a nonzero Laurent monomial with coefficient nonzero modulo three. At the parameter endpoints, we use polynomial invariants of \eqref{eq:action}. For \(H\) a set of four term indices, take in each mode the coefficient of the square-free degree-four monomial in \(\det(\sum_{s\in H}x_s X_{s,m})\) and multiply the three coefficients. The resulting polynomial is invariant because basis determinants and layer rescalings cancel. For \(H=\{4,25,28,40\}\) and \(H=\{12,15,32,38\}\) the two invariant values are \(2t^{-2}\) and \(512t^2\). Their poles prevent a finite affine limit at zero or infinity.

For a possible limit with \(t\) a unit in a discrete valuation ring \(R\), consider at each vertex the cyclic loop matrices \(A_sB_sC_s\) and their cyclic rotations. They conjugate by \(G_m\) and are invariant under term rescalings. Fifteen such matrices and the product of the first two loops give a basis of \(\operatorname{Mat}_4(R)\) at each vertex. The corresponding determinants are \(2^{-4},-2^{-4},-4t\). If the moved factors are integral, the conjugated loops are integral, so \(G_m^{-1}R^4\) is stable under all matrix units. Such a lattice is homothetic to \(R^4\), as the matrix units show by extracting and comparing the four coordinate ideals. Thus, after absorbing scalar homotheties into term rescalings, the \(G_m\) are invertible over \(R\). The unit trace makes every factor primitive, so the three remaining rescalings on each term must be integral. Their product is a unit, and hence all three are units. The residual point is therefore an ordinary specialization up to symmetry. The same argument applies to any hypothetical boundary point by selecting a valuation arc through it, passing to a finite extension of the fraction field to realize parameters and basis changes when necessary. Finite permutations preserve the two excluded conditions.
\end{proof}

\subsection*{Reconstructed curves on twelve-summand supports}
Set
\[
\begin{aligned}
 I_1&=\{4,17,34,40\}, & I_2&=\{9,25,28,35\},\\
 I_3&=\{12,19,22,32\}, & I_4&=\{15,38,44,48\}.
\end{aligned}
\]
These are the core groups used by Agarwal~\cite{agarwal}. On each union of three groups, we fix the remaining \(36\) summands of \(D(2)\) and solve the restricted equations in a gauge that fixes two nonzero factor coefficients per term. The rational reconstruction is applied to the resulting univariate series over \(\Q\). For example, in the first chart a recovered coordinate \(Y\) and a slice coordinate \(T\) satisfy
\[
 T^2Y^2+8T^2Y-512TY-4096T-2048Y=0.
\]
Passing to a rational parameter \(h\) yields
\(Y=8((1+h)^2-1)\) and \(T=-16h/((h+1)(h+2))\).
The certificate evaluates the reconstructed coordinates independently of the series solver. It checks all \(4,096\) resulting polynomial identities by degree-bounded exact interpolation. It also verifies agreement with \(D(2)\) at \(h=0\).

Write \(S_j\) for the union of three index sets in chart \(j\). Let \(Y_{s,m}^{(j)}(h)\) be the reconstructed factor for \(s\in S_j\) and mode \(m\). The data store these coordinate polynomials with a common denominator \(d_j(h)\). The certificate leaves the complementary terms at \(D(2)\), clears \(d_j(h)^3\), and verifies each of the \(4096\) polynomial identities by evaluation at one more point than the following degree bound.
\begin{center}
\small
\begin{tabular}{@{}clc@{}}
\toprule
Chart & Moving summands & Residual degree bound\\
\midrule
\(1\) & \(I_1\cup I_2\cup I_3\) & \(15\)\\
\(2\) & \(I_1\cup I_2\cup I_4\) & \(39\)\\
\(3\) & \(I_1\cup I_3\cup I_4\) & \(33\)\\
\(4\) & \(I_2\cup I_3\cup I_4\) & \(15\)\\
\bottomrule
\end{tabular}
\end{center}

To transport a chart, form in each core mode a \(16\times16\) frame \(M_m(t)\) whose first eight columns are the published basis factors in that mode and whose final eight columns are coordinate vectors completing the frame. Agarwal's normal form~\cite{agarwal} gives explicit Laurent maps on the three modes. In the second mode the map from the \(t=2\) frame to the moving frame is \(M_1(t)M_1(2)^{-1}\). The third mode has the same form with a diagonal weight inserted, and the first mode is constant. More precisely, let \(\eta=t/2\) and put third-mode weight \(\eta\) on the indices in \(I_4\), with weight \(1\) elsewhere. The third-mode map is
\[
 M_2(t)\operatorname{diag}(w_b^{-1}:b\text{ in the third-mode basis},1,\ldots,1)M_2(2)^{-1}.
\]
The three maps send the \(s\)-th core tensor summand at \(2\) to its value at \(t\) divided by \(w_s\). On each chart there is a further diagonal map on the two four-dimensional blocks in the second and third chart modes. If the three groups in that chart carry weights \(f,b,u\) in the fixed, bottom and arm positions of Agarwal's decomposition, the two maps multiply these blocks by \((f,b)\) and \((1,u/b)\), respectively. Their tensor product restores exactly the weight \(w_s\) on every selected summand. The verifier checks this termwise equality and the two frame minors as Laurent monomials. Thus these transport maps and their inverses are regular for \(t\ne0\).

Denote the three composite maps by \(L_m^{(j)}(t)\). Apply \(L_m^{(j)}(t)\) to \(Y_{s,m}^{(j)}(h)\) on the support \(S_j\), and use the original \(D_s(t)\) elsewhere. At \(h=0\) the transported partial sum is the original partial sum of \(D(t)\), and its value is independent of \(h\) by the exact chart identity. The resulting \((t,h)\) is therefore a rational two-parameter solution.

In each chart, differentiate the displayed fraction \(N(h)/d_j(h)\) at \(h=0\) and apply the transport maps with the per-term rescalings that align the base factors. The span of the infinitesimal symmetry vectors has rank \(141\) modulo three. Including the \(t\)-derivative and the four \(h\)-derivatives gives rank \(146\), both at \(t=1/4\) and at \(t=2\). These modular pivots also give a characteristic-zero lower bound. The dimension of the orbit differential in characteristic zero is at most \(141\) because of the three scalar kernel directions, so all five directions remain independent modulo the orbit over \(\Q\). Each individual chart has two independent tangent directions modulo the orbit at those points.

\begin{theorem}\label{thm:surfaces}
For each of the four reconstructed rational surfaces, the affine closure of its symmetry saturation over \(\overline{\Q}\) contains neither a zero summand nor a two-factor collision.
\end{theorem}
\begin{proof}
We outline the exact tests in the artifact and then show that they cover every affine limit. At \(t=2\), the checker computes the gcd of the vector coordinates and, for each of the \(1,128\) pairs, the gcd of the wedge coordinates in each mode. It checks that the gcd for any two-mode collision and the gcd for each zero factor have no roots away from the common denominator. At each irreducible factor of the denominator and at infinity, take the leading vector in each mode. The sum of the three factor valuations is the valuation of the summand tensor. A negative sum is a pole. If all sums are nonnegative, the certificate checks that all are zero and that the leading factors still have no two-mode collision. For a finite denominator factor these assertions are evaluated in its residue number field.

The \(36\) terms retained from \(D(t)\) determine polynomial invariants independently of the new parameter \(h\). In the third vertex, the traces of the products of cyclic loops numbered \((1,27)\) and \((3,6)\) are \(-1/(4t)\) and \(4t\) respectively. These forbid a finite affine limit at \(t=0\) or \(t=\infty\) after any gauge. If \(t\) is a unit at a valuation arc, the checker verifies that loops of the fixed terms and a product of two such loops give a Laurent-unit matrix-algebra basis at every vertex. The lattice argument from Theorem~\ref{thm:lwh} makes every vertex gauge integral modulo scalar homothety. A pole of any summand tensor in \(h\) then remains a pole in the gauged solution. A finite nonzero term has the projective leading factors checked above, transported by matrices invertible at the residual \(t\).

It remains to exclude collisions involving terms whose transport differs. In each mode all moving factors stay in the transported linear span of the chosen twelve factors at \(t=2\). For the omitted core group, the fixed factors agree projectively with transport of their \(t=2\) factors. The gcd tests at \(t=2\) therefore also cover pairs of core terms. Each term outside the sixteen-term core has a factor outside the moving span in at least two modes at every \(t\ne0\). A covector annihilating that span has Laurent-unit evaluation on the outside factor in each certified mode. Hence a moving term and an outside term can share at most one projective factor. The fixed-fixed pairs are covered by Theorem~\ref{thm:lwh}. These cases exhaust every valuation arc with generic point in the saturation and therefore every point in the affine closure.
\end{proof}


\subsection*{Minimum external support for an additional tangent}
Let \(p=D(2)\) over \(\F_3\), in the ordering of~\cite[Appendix A]{lwh}. Let \(C=I_1\cup I_2\cup I_3\cup I_4\) be the core, and let \(E=\{1,\ldots,48\}\setminus C\). Put \(K_C=\{v\in K_p:v_s=0\ (s\in E)\}\) and \(d=D'(2)\). Define the external support of \(v\) as \(\{s\in E:v_s\ne0\}\), where \(v_s\) is the triple of variations in summand \(s\). All core summands may vary.

\begin{theorem}\label{thm:external}
Among \(v\in K_p\) with \(v\notin O_p+K_C+\langle d\rangle\), the minimum external support size is \(16\). Thus the minimum is taken over representatives of tangent classes, allowing global basis changes and term rescalings.
\end{theorem}

\begin{proof}
Let \(R\subset O_p\) consist of the termwise scaling directions, and set \(W=K_p/R\). All factors are nonzero and each summand has two independent scaling directions. The certified Jacobian rank \(2153\) gives \(\dim K_p=151\), and \(\dim R=96\) gives \(\dim W=55\). The image of \(K_C\) in \(W\) has dimension four. Moreover, \(O_p+K_C\) has dimension \(145\), and \(O_p+K_C+\langle d\rangle\) has dimension \(146\). This last space represents the five known transverse directions, since the four chart derivatives are supported in \(C\) and are independent together with \(d\) modulo \(O_p\).

For \(F\subseteq E\), restrict to the summands in \(F\) and quotient by the two scaling directions \(R_s\) at each \(s\). This defines
\[
 \rho_F:W\longrightarrow\bigoplus_{s\in F}\F_3^{48}/R_s.
\]
The image of \(K_C\) in \(W\) belongs to \(\ker\rho_F\). Thus a restriction of rank \(51\) has precisely that four-dimensional kernel. Use the following partition of \(E\), with one-based indices.
\[
\begin{array}{c|rrrrrrrr}
 1&1&5&6&7&10&14&31&36\\
 2&16&20&21&23&27&29&43&46\\
 3&2&8&13&18&26&30&33&39\\
 4&3&11&24&37&41&42&45&47
\end{array}
\]
In each row there are \(40\) four-element subsets giving rank \(51\). Every one of the \(56\) five-element subsets of the row contains such a quartet. The checker regenerates the quotient basis, computes the quartet ranks and checks the cover. If \(v\) moves at most fifteen external summands, then at least seventeen are stationary. Some row therefore contains at least five stationary summands. On a certified quartet in that row the class of \(v\) lies in \(\ker\rho_F\), so \(v\in K_C+R\subseteq O_p+K_C+\langle d\rangle\).

For attainment, allow the core and the external set
\[
 S=\{5,6,8,10,11,13,16,18,20,23,24,26,36,37,41,46\}
\]
to vary. The restricted Jacobian has rank \(1467\) in \(1536\) variables, giving a \(69\)-dimensional supported kernel. Adjoining that kernel to \(O_p+K_C+\langle d\rangle\) raises the span rank from \(146\) to \(147\). The replay outputs a sparse representative and directly verifies its support and the equation \(J_pv=0\).
\end{proof}

The threshold is first-order and specific to the reduction of \(D(2)\) and the designated core. It controls tangent representatives leaving the known five-dimensional subspace. Agarwal's replacement radius measures changes to complete rank-one summands in exact decompositions~\cite{agarwal}.

\subsection*{Mixed terms and residual tensors}
For two exact deformations of the same base point \(P\), the trilinear Brent expression shows that their additive superposition is exact precisely when the mixed residual vanishes. With three or more parameters, pairwise checks suffice if at most two deformations modify each summand, because the cubic term involving three distinct parameters is then absent. At \(D(2)\) we checked the six pairs of stored core curves by exact rational evaluation. At parameter \(1\) on both curves, the two individual identities hold, while the additive superposition fails for every pair. Thus additive superposition in the stored factor coordinates cannot produce a joint family of these curves. A second test rules out a more general sequential strategy. At the parameter value (1) on each chart, for each of the three other stored supports, the restricted Jacobian has rank (552) modulo (101). A nonzero minor modulo (101) gives the same lower bound over (Q), and (24) independent rescalings give the matching upper bound in (576) variables. Hence at a generic point on each chart the restricted Brent scheme has tangent dimension \(24\). The \(24\)-dimensional scaling orbit through that point supplies the matching local dimension lower bound. The point is smooth in the restricted scheme, and its scaling orbit is locally open. In a regular joint parametrization with a second parameter supported on one of those sets, an independent derivative at (D(2)) would remain independent on a nonempty open part of the first curve, contradicting this rank. A joint construction may move additional terms. The overlapping supports also require accounting for triple mixed terms when testing more than two parameters.

A further boundary question is whether a block of divergent summands can be replaced while the remaining finite summand limits are held fixed. After subtracting the finite limits from the multiplication tensor, a rank bound on the residual tensor can rule out such a replacement. Flattenings give one test. If a residual has an \(r\)-dimensional matrix slice space and a \(2\times2\) minor on that space is a nonzero multiple of \(\ell^2\) for a nonzero linear form \(\ell\), then its rank is at least \(r+1\). Indeed, a length-\(r\) decomposition would make the slice space spanned by \(r\) rank-one matrices, all of which must lie in \(\ker\ell\). This test arose while studying singular limits in deformations of the 93-term \(5\times5\) scheme of Moosbauer and Poole~\cite{mp}. A block-replacement analysis for DPS48 would require its own residual tensors.

\paragraph{Rank and collision scan.} The frozen catalogue of \(21\) displayed rational tables has exact mod-three Jacobian rank \(2,153\) and orbit differential rank \(141\) in each case. It includes the PLinOpt natural-basis tables and composed encodings, the Moran--Schwartz--Yuan table, and Agarwal's twelve-term replacement. We also audit the good reductions on a finite rational grid on the reconstructed curves. The discrete walk and the grid yielded no exact 47-term rational identity. The latter is a search result, and the boundary theorem concerns the four certified surfaces.

\section{Comparison with prior results}
\begin{center}
\small
\begin{tabularx}{\textwidth}{@{}p{.24\textwidth}>{\raggedright\arraybackslash}X>{\raggedright\arraybackslash}X@{}}
\toprule
Work & Object audited & Scope of conclusion\\
\midrule
Tomar~\cite{tomar} & Five binary 47-term representatives & Bounded complete-summand replacement.\\
Kauers--Moosbauer~\cite{km} & Binary flip-graph tables & First Hensel failures at scale.\\
Li--Wang--Hu~\cite{lwh} & DPS48 & Tangent and orbit computation, deflation evidence, and exact one-parameter family.\\
Agarwal~\cite{agarwal} & \(D(t)\) and a twelve-term replacement & Bounded-subset rigidity and exact transport of completion charts.\\
This work & Binary points and DPS48 & Exact local geometry, ramified obstruction, an orbit certificate, reconstructed rational subfamilies, scoped boundary exclusions, and an exact transverse support threshold.\\
\bottomrule
\end{tabularx}
\end{center}

\section{Reproducibility}
The public repository contains the manuscript, factor tables, certificates, and an offline verifier. Use the files from the version tag for this manuscript.
\begin{quote}
\url{https://github.com/edsonphx/brent-47-48-schemes}\\
\url{https://github.com/edsonphx/brent-47-48-schemes/tree/v1.1.0}
\end{quote}
Run \texttt{bash reproduce.sh} with Python, NumPy and SymPy from the accompanying archive. The scripts \texttt{equivalence.py}, \texttt{hensel.py}, \texttt{families.py} and \texttt{dps48\_support.py} also run independently. The replay reads frozen coefficient tables, reevaluates all tensor identities, checks the parity witnesses and orbit certificates, and verifies the rational polynomial identities, boundary certificates, and external-support threshold. It also emits the exact pairwise superposition counterexamples. Source revisions and reduction conventions are recorded in the manifest and data notes. The numerical collision search is exploratory and is described separately from the exact evidence.

\subsection*{AI-assisted research statement}
An OpenAI model assisted with literature review, exploratory calculations, certificate code, and drafting. The computational statements in this manuscript are presented with replayable checks and explicit mathematical arguments.

\appendix
\section{Binary cokernel witnesses}
\begin{center}
\begin{tabular}{@{}lrr@{}}
\toprule
Input & Projection weight & Dual witness weight\\
\midrule
AT47 &12&520\\
KM47 &149&290\\
C3 &240&606\\
C4 &532&620\\
C5 &32&253\\
Gamma47-B &288&682\\
\bottomrule
\end{tabular}
\end{center}
A projection weight counts nonzero coordinates after eliminating the image pivots from \(F_{47}(p)/2\). It depends on the displayed coefficient table and elimination order. A dual witness weight counts the equations appearing in its parity functional.

\begin{thebibliography}{9}
\bibitem{strassensquare} E. Soares de Souza, \emph{Exact local dimension of the Brent variety \(V(4,4,4\mid49)\) at the recursive Strassen-square point via quadratic obstructions}, 2026. \url{https://doi.org/10.5281/zenodo.24449376}
\bibitem{tomar} T.~S. Tomar, \emph{Local optimality of rank-47 schemes for \(4\times4\) matrix multiplication over \(\F_2\)}, 2026. \url{https://www.mefyi.com/papers/rank-47-local-optimality/}
\bibitem{km} M. Kauers and J. Moosbauer, \emph{Flip graphs for matrix multiplication}, ISSAC 2023. \url{https://arxiv.org/abs/2212.01175}
\bibitem{lwh} X. Li, Y. Wang, and S. Hu, \emph{Substitution and quotient of the isotropy group action}, 2026. \url{https://arxiv.org/abs/2607.15069}
\bibitem{dps} J.-G. Dumas, C. Pernet, and A. Sedoglavic, \emph{A non-commutative algorithm for multiplying \(4\times4\) matrices using \(48\) non-complex multiplications}, 2025. \url{https://arxiv.org/abs/2506.13242}
\bibitem{mp} J. Moosbauer and M. Poole, \emph{Flip graphs with symmetry and new matrix multiplication schemes}, ISSAC 2025. \url{https://arxiv.org/abs/2502.04514}
\bibitem{agarwal} A. Agarwal, \emph{Certified local rank and uniqueness barriers for a 48-term matrix-multiplication decomposition}, 2026. \url{https://arxiv.org/abs/2609.13596}
\end{thebibliography}
\end{document}
